\section{Описание стенда и подготовка}

\subsection{Конфигурация виртуального стенда}

Для выполнения работы использовались две виртуальные машины (VirtualBox):
\begin{itemize}
  \item \textbf{gateway}~--- Ubuntu~20.04.5~LTS Server, 2 сетевых адаптера:
    \begin{itemize}
      \item \texttt{enp0s3} (NAT): \texttt{10.0.2.15/24}, шлюз \texttt{10.0.2.2};
      \item \texttt{enp0s8} (Внутренняя сеть): \texttt{192.168.\cfgLabStudentN.1/24}.
    \end{itemize}
  \item \textbf{Desktop1}~--- Ubuntu~Desktop, 1 адаптер (внутренняя сеть).
\end{itemize}

\subsection{Настройка netplan на сервере}

\begin{lstlisting}
sudo su
nano /etc/netplan/00-installer-config.yaml
\end{lstlisting}

\begin{lstlisting}
network:
  version: 2
  renderer: networkd
  ethernets:
    enp0s3:
      dhcp4: no
      addresses:
        - 10.0.2.15/24
      gateway4: 10.0.2.2
      nameservers:
        addresses: [8.8.8.8]
    enp0s8:
      dhcp4: no
      addresses:
        - 192.168.N.1/24
\end{lstlisting}

\begin{lstlisting}
netplan apply
ip a
ping ya.ru
\end{lstlisting}

После \texttt{netplan apply} интерфейс \texttt{enp0s3} получил адрес
\texttt{10.0.2.15/24}, \texttt{enp0s8}~--- \texttt{192.168.\cfgLabStudentN.1/24}.
Доступ в Интернет с сервера подтверждён \texttt{ping~ya.ru}.

\subsection{Настройка клиента Desktop1 (вручную)}

До развёртывания DHCP-сервера IP-адрес клиенту назначался вручную через графический
интерфейс \texttt{nm-connection-editor}:
\begin{itemize}
  \item IP: \texttt{192.168.\cfgLabStudentN.10};
  \item Маска: \texttt{255.255.255.0};
  \item Шлюз: \texttt{192.168.\cfgLabStudentN.1};
  \item DNS: \texttt{192.168.\cfgLabStudentN.1}.
\end{itemize}

Проверка связи: \texttt{ping~192.168.\cfgLabStudentN.1}~--- запросы проходят.
